% GENERATED FILE. Edit canonical lessons, then run npm run generate:books.
% canonical-source-hash: fnv1a64:9c70ed855734e972
% canonical-lessons: UR-C324-sunana, UR-C324-dastan, UR-C324-qissa, UR-C324-paheli, UR-C324-kahavat

\chapter{A Tale of Long Ago}
\label{ch:ur-a-tale-of-long-ago}
\hlchaptermodality{324}

\begin{chapteropening}
\textbf{By the end of this chapter:} \emph{I can begin and tell a traditional tale.}

Complete the last lesson of chapter 324: I can begin and tell a traditional tale.
\end{chapteropening}

\section[\texorpdfstring{\ur{سنانا} (sunānā)}{سنانا (sunānā)}]{\ur{سنانا} (sunānā) — to tell a story; to recite}

\label{lesson:UR-C324-sunana}



\begin{quote}
Before the new one: say the Urdu for to accept, then the Urdu for to keep up.
\end{quote}

\subsection*{You'll want to know: \ur{سنانا}}
\textbf{\ur{سنانا}} — \emph{sunānā} — \textquotedblleft{}to tell a story; to recite\textquotedblright{}.

\begin{quote}
\ur{امی} \ur{نے} \ur{ہمیں} \ur{ایک} \ur{گیت} \ur{سنایا۔}
\end{quote}

\begin{quote}
\emph{ammī ne hameṅ ek gīt sunāyā.}
\end{quote}

\begin{quote}
Mum sang us a song.
\end{quote}

Past, with \ur{نے} (he or she …): \ur{اس} \ur{نے} \ur{نظم} \ur{سنائی}, \emph{us ne nazm sunāī}.

\begin{cousinweb}
is the causing partner of \textbf{\ur{سننا}}, to hear: CDIAL lists Hindi \emph{sunānā}, \textquotedblleft{}informs\textquotedblright{}, under Sanskrit \emph{śṛṇoti}, \textquotedblleft{}hears\textquotedblright{} (CDIAL 12598). To \textbf{\ur{سنانا}} is to make someone hear: a story, a song, a poem.
\end{cousinweb}

\begin{grammarlens}[title={What you've built}]
One more word for telling what happened.
\end{grammarlens}

\subsection*{Your turn}
\begin{itemize}
\raggedright
  \item \emph{Say it:} \emph{sunānā}
  \item \emph{Say it:} \emph{sunānā}, once more
  \item \emph{Say it:} \emph{ammī ne hameṅ ek gīt sunāyā.} — Mum sang us a song.
  \item \emph{Recall:} say the Urdu for to accept, then the Urdu for to keep up, then say \emph{sunānā} again
\end{itemize}

\begin{tcolorbox}[breakable,colback=teal!4,colframe=teal!35!black,title={Before you move on}]
What does \ur{سنانا} mean? (\textquotedblleft{}To tell a story; to recite\textquotedblright{}.) Say it once more.
\end{tcolorbox}
\section[\texorpdfstring{\ur{داستان} (dāstān)}{داستان (dāstān)}]{\ur{داستان} (dāstān) — an epic tale; a long story}

\label{lesson:UR-C324-dastan}



\begin{quote}
Before the new one: say the Urdu for to keep up, then the Urdu for to tell a story.
\end{quote}

\subsection*{You'll want to know: \ur{داستان}}
\textbf{\ur{داستان}} — \emph{dāstān} — \textquotedblleft{}an epic tale; a long story\textquotedblright{}.

\begin{quote}
\ur{دادا} \ur{نے} \ur{ایک} \ur{لمبی} \ur{داستان} \ur{سنائی۔}
\end{quote}

\begin{quote}
\emph{dādā ne ek lambī dāstān sunāī.}
\end{quote}

\begin{quote}
Grandfather told a long tale.
\end{quote}

Feminine.

\begin{cousinweb}
is usually traced to Persian \emph{dāstān}, \textquotedblleft{}a tale\textquotedblright{}. Urdu's old \textbf{\ur{داستان}} went on for nights, with kings, giants and fairies.
\end{cousinweb}

\begin{grammarlens}[title={What you've built}]
One more word for telling what happened.
\end{grammarlens}

\subsection*{Your turn}
\begin{itemize}
\raggedright
  \item \emph{Say it:} \emph{dāstān}
  \item \emph{Say it:} \emph{dāstān}, once more
  \item \emph{Say it:} \emph{dādā ne ek lambī dāstān sunāī.} — Grandfather told a long tale.
  \item \emph{Recall:} say the Urdu for to keep up, then the Urdu for to tell a story, then say \emph{dāstān} again
\end{itemize}

\begin{tcolorbox}[breakable,colback=teal!4,colframe=teal!35!black,title={Before you move on}]
What does \ur{داستان} mean? (\textquotedblleft{}An epic tale; a long story\textquotedblright{}.) Say it once more.
\end{tcolorbox}
\section[\texorpdfstring{\ur{قصہ} (qissā)}{قصہ (qissā)}]{\ur{قصہ} (qissā) — a tale; an anecdote}

\label{lesson:UR-C324-qissa}



\begin{quote}
Before the new one: say the Urdu for to tell a story, then the Urdu for an epic tale.
\end{quote}

\subsection*{You'll want to know: \ur{قصہ}}
\textbf{\ur{قصہ}} — \emph{qissā} — \textquotedblleft{}a tale; an anecdote\textquotedblright{}.

\begin{quote}
\ur{یہ} \ur{قصہ} \ur{سب} \ur{جانتے} \ur{ہیں۔}
\end{quote}

\begin{quote}
\emph{yih qissā sab jānte haiṅ.}
\end{quote}

\begin{quote}
Everyone knows this story.
\end{quote}

Masculine.

\begin{cousinweb}
is usually traced, through Persian, to Arabic \emph{qiṣṣa}, \textquotedblleft{}a story\textquotedblright{}. A \textbf{\ur{قصہ}} is short; a \textbf{\ur{داستان}} is long.
\end{cousinweb}

\begin{grammarlens}[title={What you've built}]
One more word for telling what happened.
\end{grammarlens}

\subsection*{Your turn}
\begin{itemize}
\raggedright
  \item \emph{Say it:} \emph{qissā}
  \item \emph{Say it:} \emph{qissā}, once more
  \item \emph{Say it:} \emph{yih qissā sab jānte haiṅ.} — Everyone knows this story.
  \item \emph{Recall:} say the Urdu for to tell a story, then the Urdu for an epic tale, then say \emph{qissā} again
\end{itemize}

\begin{tcolorbox}[breakable,colback=teal!4,colframe=teal!35!black,title={Before you move on}]
What does \ur{قصہ} mean? (\textquotedblleft{}A tale; an anecdote\textquotedblright{}.) Say it once more.
\end{tcolorbox}
\section[\texorpdfstring{\ur{پہیلی} (pahelī)}{پہیلی (pahelī)}]{\ur{پہیلی} (pahelī) — a riddle}

\label{lesson:UR-C324-paheli}



\begin{quote}
Before the new one: say the Urdu for an epic tale, then the Urdu for an anecdote.
\end{quote}

\subsection*{You'll want to know: \ur{پہیلی}}
\textbf{\ur{پہیلی}} — \emph{pahelī} — \textquotedblleft{}a riddle\textquotedblright{}.

\begin{quote}
\ur{کوئی} \ur{پہیلی} \ur{کا} \ur{جواب} \ur{نہیں} \ur{جانتا} \ur{تھا۔}
\end{quote}

\begin{quote}
\emph{koī pahelī kā javāb nahīṅ jāntā thā.}
\end{quote}

\begin{quote}
Nobody knew the answer to the riddle.
\end{quote}

Feminine.

\begin{cousinweb}
\textbf{\ur{پہیلی}} (\emph{pahelī}) is inherited from Sanskrit \emph{prahēlikā}, \textquotedblleft{}a riddle\textquotedblright{} (CDIAL 8911).
\end{cousinweb}

\begin{grammarlens}[title={What you've built}]
One more word for telling what happened.
\end{grammarlens}

\subsection*{Your turn}
\begin{itemize}
\raggedright
  \item \emph{Say it:} \emph{pahelī}
  \item \emph{Say it:} \emph{pahelī}, once more
  \item \emph{Say it:} \emph{koī pahelī kā javāb nahīṅ jāntā thā.} — Nobody knew the answer to the riddle.
  \item \emph{Recall:} say the Urdu for an epic tale, then the Urdu for an anecdote, then say \emph{pahelī} again
\end{itemize}

\begin{tcolorbox}[breakable,colback=teal!4,colframe=teal!35!black,title={Before you move on}]
What does \ur{پہیلی} mean? (\textquotedblleft{}A riddle\textquotedblright{}.) Say it once more.
\end{tcolorbox}
\section[\texorpdfstring{\ur{کہاوت} (kahāvat)}{کہاوت (kahāvat)}]{\ur{کہاوت} (kahāvat) — a saying; a proverb}

\label{lesson:UR-C324-kahavat}



\begin{quote}
Before the new one: say the Urdu for an anecdote, then the Urdu for a riddle.
\end{quote}

\subsection*{You'll want to know: \ur{کہاوت}}
\textbf{\ur{کہاوت}} — \emph{kahāvat} — \textquotedblleft{}a saying; a proverb\textquotedblright{}.

\begin{quote}
\ur{دادی} \ur{ہر} \ur{بات} \ur{پر} \ur{کہاوت} \ur{سناتی} \ur{تھیں۔}
\end{quote}

\begin{quote}
\emph{dādī har bāt par kahāvat sunātī thīṅ.}
\end{quote}

\begin{quote}
Grandmother had a proverb for everything.
\end{quote}

Feminine.

\begin{cousinweb}
\textbf{\ur{کہاوت}} (\emph{kahāvat}) is inherited: CDIAL lists the Hindi cousin \emph{kahāwat}, \textquotedblleft{}a saying\textquotedblright{}, under a form it reconstructs (CDIAL 2706). It belongs with \textbf{\ur{کہنا}}, to say.
\end{cousinweb}

\begin{grammarlens}[title={What you've built}]
One more word for telling what happened.
\end{grammarlens}

\subsection*{Your turn}
\begin{itemize}
\raggedright
  \item \emph{Say it:} \emph{kahāvat}
  \item \emph{Say it:} \emph{kahāvat}, once more
  \item \emph{Say it:} \emph{dādī har bāt par kahāvat sunātī thīṅ.} — Grandmother had a proverb for everything.
  \item \emph{Recall:} say the Urdu for an anecdote, then the Urdu for a riddle, then say \emph{kahāvat} again
\end{itemize}

\begin{tcolorbox}[breakable,colback=teal!4,colframe=teal!35!black,title={Before you move on}]
What does \ur{کہاوت} mean? (\textquotedblleft{}A saying; a proverb\textquotedblright{}.) Say it once more.
\end{tcolorbox}
